Redis Cluster在多个主节点间分配键空间，每个主节点再配副本。键先映射到16384个槽，槽再映射到当前主节点。节点位置变化只改变后一层关系，键的内容和业务身份不必随地址变化。客户端缓存槽表，命令直接到当前主节点，收到重定向时再刷新~\cite{redis-cluster-spec}。

\subsection{键、槽与同一执行边界}

一般键按CRC16计算并取模16384。哈希标签提供受控的共同归属：找到键中第一个左花括号及其后第一个右花括号，二者之间非空时，用其中内容作为散列输入，否则使用整个键。下面继续检验把业务权威直接放进Redis的假设：\code{accepted:\{b1\}:S1}、\code{pending:\{b1\}:C1}和\code{session:\{b1\}:E1}落在同一槽，其中S1只是完整四分量S在键名中的编码代号。共同标签让多键事务、脚本与多键命令有机会在同一节点执行，也把这些键的容量和热点放在同一分片。

\CEighteenFigure{ch18-15}{完整状态的同槽组织与热点范围}{fig:18-hash-tag}

标签的业务范围由设计者选择。按每个S打标签，只保证这个S相关键共同归属，同一会话中的准入状态、跨S共享的退款额度或库存对象可能落到别处。按整个b1打标签能扩大共同执行边界，也会让b1全部压力落在一个槽。把所有账户都放进同一个标签，会把本来可以分散的压力集中到一个主节点。

Cluster通常会拒绝跨槽的多键访问，客户端得到\code{CROSSSLOT}。需要同一原子状态迁移的应用，先检查数据归属与命令所支持的边界。组合商品的两个库存对象如果分属不同分区，就由第19章的参与者准备和持久决定连接起来。哈希标签选择数据局部性，跨分区事务负责跨局部性的业务条件。

对错误节点发请求，\code{MOVED}告诉客户端该槽当前归属的节点，客户端据此更新长期路由。读取副本另有只读连接和可能陈旧的语义，写请求按槽所有者及当前角色处理。槽表和节点状态经集群总线传播，在传播尚未完成时，不同节点可能持有不同视图，客户端通过重定向逐步找到服务节点。

\subsection{传统迁槽的逐键过程}

传统协议把源节点A的槽设为\code{MIGRATING B}，把目标B对应槽设为\code{IMPORTING A}。正式归属暂时指向A。A对存在的键继续服务，对已经不在A的键返回\code{ASK}，要求这一次请求到B。B收到普通请求时通常按现有归属返回MOVED；客户端先在同一目标连接发送\code{ASKING}，再发送刚才那条请求，B才按迁入状态处理。

\CEighteenFigure{ch18-16}{传统迁槽中的现存键、ASK与一次性ASKING}{fig:18-legacy-migration}

搬迁器用\code{MIGRATE}把指定键序列化并送到B，收到成功后从A删除相应键。迁移单个键期间，命令处理与传输形成阻塞成本，大键尤其容易拉长延迟。网络I/O超时可能让调用方无法立即分辨目标是否已接收，迁移工具随后核对原键与目标键的实际状态~\cite{redis-cluster-spec}。

同槽多键操作在迁槽中也可能遇到键分散。设\code{accepted:\allowbreak\{b1\}:S1}已经迁到B，\code{pending:\allowbreak\{b1\}:C1}还在A。A收到同时访问两键的请求时，既有现存键又有缺失键，无法在本地执行完整操作，也无法把请求整体转到B，于是返回\code{TRYAGAIN}。目标在迁入状态下遇到不完整的多键集合，也可能采用相应拒绝。客户端退避后用原输入重试，搬迁器继续推进完整槽，缺键在这里反映迁槽阶段。

全部键搬完后，管理者把槽的正式归属设为B，并传播给其他节点。A此后对旧路由返回MOVED，客户端把后续请求长期改发B。ASK授权一次临时转向，MOVED表达当前正式路由，两者在客户端的缓存动作不同。搬迁器退出时，MIGRATING/IMPORTING与实际键位置可能保留，恢复工具从双方状态继续完成数据和归属交接。

ASKING的连接范围还会影响客户端实现。客户端在到B的连接c1发送ASKING，却把下一条GET交给连接池中的c2，c2没有这次标记，B按普通路由处理。正确路径把ASKING和被重定向命令留在同一目标连接，并保留原命令参数。若客户端在收到ASK时就永久把整槽改为B，尚在A的其他键又会触发重定向，增加往返。

\CEighteenFigure{ch18-26-redirect}{同槽双键迁移中的请求分叉与连接范围}{fig:18-redirect-branches}

\subsection{新槽迁移的快照、流与接管}

从Redis 8.4开始，Cluster提供原子槽迁移命令。在8.10.2中，目标节点执行以下命令启动相应槽范围的导入任务，取得任务ID后可以查询状态或取消\footnote{ASM是Atomic Slot Migration的简称，指该版的槽快照、增量、排空与接管协议。传统SETSLOT加逐键MIGRATE的路径在8.10.2中与它并存。}。这里的原子性指向该协议所组织的槽迁移和接管，故障条件要沿其实际状态路径判断~\cite{redis8102-asm}。

\begin{BookCode}
CLUSTER MIGRATION IMPORT start end
\end{BookCode}

新协议让源节点生成所迁槽的快照，目标接收快照的同时缓冲后续修改流。源侧继续处理相应请求，写入变更随迁移流发往目标。目标加载快照并消费积累的修改，当落后缩小到适合交接时，源侧暂停相应写入，让目标排空剩余流，随后执行接管与配置传播。源侧观察到新归属后完成任务并清理旧槽数据。

\CEighteenFigure{ch18-17}{ASM的快照、增量、排空与接管}{fig:18-asm}

这条路径把大部分数据传输移到后台，使最终写暂停集中在排空与归属交接附近。代价转移为快照成本、流缓冲、网络带宽和目标应用速度。写入生成速度长期高于目标消费速度时，后缀不能自然缩小，任务进入受控限流、暂停或失败处理。状态接口中的阶段、重试、最后错误及写暂停时间，分别帮助观察任务推进到哪里。

数据完整性还取决于应用放入这个槽的对象。ASM迁移槽内状态，独立平台的X1和另一个槽中的T1各在自己的保存域。标签与业务结构在迁移前把原身份和恢复责任闭合，槽迁移才有可能把这些本地事实整体交接。第17章业务分区迁移还显式核对冻结图像、目录CAS和目标激活，因为那里迁移的是完整业务权威及其执行资格。

当前源码对接管超时给出了一条具体故障历史。源节点已经暂停写入并等待目标发布归属，目标却暂时失联。为避免源节点永远暂停，任务超时失败并恢复源写入。随后目标迟到的归属广播到达，带着提高后的\code{configEpoch}，配置最终可能以目标为准。源在恢复写入到归属收敛之间接收的写，没有继续复制到目标，因而会丢失~\cite{redis8102-asm}。

\CEighteenFigure{ch18-18}{接管超时后，源恢复写入与迟到归属发生冲突}{fig:18-asm-timeout}

迁移阶段、超时动作和新旧写资格在这条历史中共同决定恢复点。Redis槽若保存唯一受理，恢复写入期间新建的C及其待办可能随w一起消失，而已离开节点的付款请求继续生效。采用该权威方案的业务据此限定准入与恢复范围，并从独立证据找回已经发送的责任。购买系统只把可重建投影放进这些槽，归属收敛后清空受影响投影，查询从原权威取得足够新的版本。

\subsection{拓扑代次、故障检测与成员变化}

Cluster节点通过心跳和gossip交换已知节点、槽位及故障观察。某节点本地观察到超时，可以标记对端PFAIL；收集有效时间窗内多数主节点的故障报告后，形成FAIL并传播。故障观察用于触发提升，副本接替还要经过选举授权。普通数据复制沿各主副本关系进行，每条业务写入并不经过全部主节点共同提交~\cite{redis-cluster-spec}。

集群用\code{currentEpoch}为配置事件提供逻辑代次，节点看见更大值就推进自己的观察。槽归属声明带有\code{configEpoch}，随主节点所声称的槽传播，用来比较冲突配置。合格副本竞选时增加代次并请求有槽主节点投票，每个投票者在相应代次内受投票规则约束。得到多数授权后，新主以更新的配置代次接替槽。配置代次比较解决槽归属声明的先后，已回复业务写入是否存在，则由新主实际获得的数据历史决定。

提升时的多数按拥有槽的主节点计算，普通副本和没有槽的空主节点不会因加入集群就增加这一投票权。三主三副本的部署中，R1接替M1，需要满足自己的候选条件并取得主节点多数授权。M1和某个客户端被隔离在少数侧，M2、M3与R1位于另一侧时，另一侧可以推进接替，M1在检测窗口后停止相应服务。检测窗口之前M1已经确认而未复制给R1的写，可能随最终配置收敛丢失。

\CEighteenFigure{ch18-27-cluster-vote}{槽配置的多数授权与每个主节点的数据复制分别推进}{fig:18-cluster-votes}

拓扑恢复同时依赖节点配置保存。节点把学习到的代次与槽关系维护到自己的集群配置文件，启动后据此继续和其他节点交换信息。这份文件含节点自身身份、角色和已知拓扑，部署入口为每个节点分配独立数据目录和配置文件，再通过集群协议建立关系。

成员变化也沿数据路径完成。添加空节点后，先把需要分担的槽迁入，再让客户端和其他节点学习新拓扑；删除节点前，先迁出槽并安排其副本关系。节点恢复时，根据当前归属和复制关系接入，旧主分叉状态按当前主历史处理。配置管理与数据责任交接一起完成，迁出的槽才有明确的新承载。

槽的归属交接确定接下来从哪里读取键，文件和复制前缀确定这些键能恢复到哪里。更大的状态图像还会跨越多个文件或对象：只有各块齐备，并由同一份已发布清单选中，恢复者才能取得一个完整版本。下一章沿这条发布路径进入分布式存储，再展开跨区间修改怎样形成共同事务决定。
\FloatBarrier
